%% Representable free models, arbitrary generators, and carrier boundaries.

\section{Free Lawvere models and carrier conventions}

\begin{annexitem}{annex:p10.lawvere-free.1}{thm:p10-lawvere-finite-free}{Yoneda and the finite free models}
\begin{proof}
\paragraph{Representables preserve products.}
For a finite product cone in \(\mathbb T\), its universal property gives
a bijection from arrows with domain \(n\) into the product to tuples of
arrows from \(n\) into its factors.
These are exactly the canonical product comparisons for \(R_n\).
In particular,
\[
 \mathbb T(n,k)\longrightarrow\mathbb T(n,1)^k,
 \qquad a\longmapsto(p_i^k\circ a)_{i<k}
\]
is a bijection.  For \(k=0\), both sides are singletons because \(0\)
is terminal.  Postcomposition preserves identities and composition,
so \(R_n\) is a product-preserving functor, hence a model.

\paragraph{The free generators.}
Apply Yoneda to \(\mathbb T^{\mathrm{op}}\): its presheaves are
covariant functors on \(\mathbb T\), and its representable at \(n\)
is \(\mathbb T(n,-)\).  Thus
\[
 \operatorname{Nat}(R_n,M)\cong M(n),
 \qquad \alpha\longmapsto\alpha_n(1_n).
\]
The inverse sends \(z\in M(n)\) to the transformation
\[
 \alpha^z_k(a)=M(a)(z)\qquad(a:n\to k).
\]
Its naturality follows from \(M(t a)=M(t)M(a)\).
Conversely, naturality of \(\alpha\) at \(a:n\to k\) gives
\(\alpha_k(a)=M(a)(\alpha_n(1_n))\), proving both inverse identities.

Compose this bijection with \(c_n^M\).
For each \(i<n\), naturality at \(p_i^n:n\to1\) gives
\[
 M(p_i^n)(\alpha_n(1_n))=\alpha_1(p_i^n).
\]
Hence the resulting map to \(M(1)^n=\mathsf{Set}(n,M(1))\) is exactly
restriction to \(\eta^n\).
For any natural transformation \(b:M\to N\),
\[
 c_n^N\circ b_n=(b_1)^n\circ c_n^M,
\]
as follows by comparing each projection.
Together with Yoneda naturality this proves naturality of the free
bijection for all model homomorphisms.
At \(n=0\), \(M(0)\) is a singleton, so there is exactly one map
\(R_0\to M\); this proves initiality.
Its underlying carrier is its value \(R_0(1)=\mathbb T(0,1)\).

\paragraph{The opposite theory.}
For \(f:m\to n\), precomposition defines
\[
 R_n\longrightarrow R_m,
 \qquad (a:n\to k)\longmapsto a\circ f.
\]
The identity and composition laws follow from composition in \(\mathbb T\).
Yoneda, now with \(M=R_m\), identifies every transformation
\(R_n\to R_m\) with exactly one \(f\in\mathbb T(m,n)\).
This proves full faithfulness and the asserted isomorphism when the
chosen representables are indexed by their objects \(n\).

For a finite set \(X\), fix a bijection \(n\to X\).
Any model free on \(X\) and \(R_n\) are then free on the same indexed
generators. Their universal properties give arrows in both directions
respecting those generators. Each composite is the identity by the
same uniqueness property.
The isomorphism is unique with that compatibility.
This proves the statement about all finite free models without choosing
a free basis for every isomorphic model simultaneously.
\end{proof}
\end{annexitem}

\begin{annexitem}{annex:p10.lawvere-free.2}{thm:p10-lawvere-free-adjunction}{Terms, arbitrary free models, and the adjunction}
\begin{proof}
\paragraph{The term set and its semantic congruence.}
Use formal tagged variables for the elements of \(X\).
The term-evaluation, substitution, and homomorphism inductions of
Annex~\ref{annex:p10.eq.1} apply whenever an algebra and an assignment
are given, whether its carrier is empty or nonempty.
If the term set is empty, it has no nullary operation; every
positive-arity term operation on it is the unique empty function.
It is therefore an all-carrier algebra.

The relation \(\theta_X\) is a congruence: it is the common condition
of belonging to the kernels of all evaluation maps into all models of
\(E\).  Equivalence and operation compatibility hold in each kernel,
and hence in this common condition.
It is a relation on the set \(T_\Sigma(X)^2\); no product indexed by
the proper class of all algebras is formed.
The quotient operations are consequently well-defined.

\paragraph{The quotient satisfies every defining equation.}
First suppose \(Q_X\ne\varnothing\).
Fix an equation \(s\approx t\) of \(E\) in its finite context \(Y\)
and an assignment \(a:Y\to Q_X\).
Choose representative terms \(r_y\) for these finitely many values.
They give a substitution \(r:Y\to T_\Sigma(X)\).
Induction on terms gives
\[
 s^{Q_X}(a)=[s r],\qquad t^{Q_X}(a)=[t r].
\]
For every all-carrier model \(A\) of \(E\) and every \(v:X\to A\),
the substitution law gives
\[
 (s r)^A(v)=s^A(y\mapsto r_y^A(v))
           =t^A(y\mapsto r_y^A(v))=(t r)^A(v).
\]
The middle equality uses the equation in its declared context.
Thus \((s r,t r)\in\theta_X\), proving equality in \(Q_X\).
This proof includes \(Y=\varnothing\), when the substitutions are empty
and the terms are ground.
Only a finite family of representatives is chosen in any instance.

If \(Q_X=\varnothing\), then \(T_\Sigma(X)=\varnothing\): a quotient
of a nonempty set has an equivalence class.
Consequently \(X=\varnothing\) and \(\Sigma\) has no nullary symbols.
Every term in this signature contains a variable, so every well-formed
equation has a nonempty declared variable context.
There is no assignment from such a context to the empty carrier;
all these identities are satisfied.
This proves \(Q_X\in\operatorname{Alg}_{\rm all}(\Sigma,E)\) in every case.

\paragraph{The universal property and functoriality.}
For \(v:X\to A\), set \(\bar v([t])=t^A(v)\).
The definition of \(\theta_X\) makes the value independent of the
representative, and the evaluation recursion makes it a homomorphism.
It extends \(x\mapsto v(x)\) along \(x\mapsto[x]\).
Every homomorphism with these generator values has the same value on
every term, by induction using preservation of all operations, including
nullaries.  Hence the extension is unique.
For an empty term set this is the unique map from the empty algebra;
its positive-arity homomorphism laws are vacuous.

For \(u:X\to Y\), the unique homomorphism
\(Q(u):Q_X\to Q_Y\) extending \(x\mapsto[u(x)]\) sends
\([t]\) to \([t[u]]\).
Here \(t[u]\) replaces each tagged generator by the tagged generator
with image label \(u(x)\).
The extension uniqueness proves \(Q(1_X)=1_{Q_X}\) and
\(Q(vu)=Q(v)Q(u)\).
Thus the construction is functorial, with no selected representatives
for elements of the quotient.

\paragraph{Transfer to models with the actual tuple comparisons.}
Use the fixed recovery presentation and \(\rho\) in the statement.
For each \(M\), pull it back along \(\rho\), and let \(A_M\) be
the recovered \((\Sigma_{\mathbb T},E_{\mathbb T})\)-algebra of
Theorem~\ref{thm:p10-lawvere-models}.
Its carrier is \(M(1)\).
The recovery isomorphism fixes objects and the structural projections,
so the comparisons \(c_n^M\) apply unchanged.

Given \(v:X\to M(1)\), the free-algebra property gives
\(\bar v:Q_X\to A_M\).
Define
\[
 \alpha_n=(c_n^M)^{-1}\circ(\bar v)^n:
                Q_X^n\longrightarrow M(n).
\]
The model--algebra correspondence, which identifies all homomorphisms
with natural transformations, proves that these components form
\(\alpha:\mathcal F_{\mathbb T}(X)\to M\).
Conversely, a natural transformation \(\alpha\) gives the algebra
homomorphism
\[
 Q_X\longrightarrow A_M,\qquad q\longmapsto\alpha_1((q)).
\]
These assignments are inverse by the same correspondence.
Since \(p_0^1=1_1\), the comparison \(c_1^M\) sends \(a\) to \((a)\).
Thus the displayed \(\alpha\) satisfies
\(\alpha_1(([x]))=\bar v([x])=v(x)\).
Restriction along \(\eta_X(x)=([x])\) is therefore the asserted
free-model bijection, with exactly the actual underlying set \(M(1)\).
At \(n=0\), the component is the unique map between the two singleton
sets; it does not assert that \(M(1)\) is nonempty.

Let \(\mathcal F_{\mathbb T}(u)\) have components \(Q(u)^n\), transported
along \(\rho^{-1}\). The algebra-to-model construction makes these
natural transformations. The proved functor laws of \(Q\) give those
of \(\mathcal F_{\mathbb T}\).
Its generator maps satisfy
\[
 U\mathcal F_{\mathbb T}(u)\circ\eta_X=\eta_Y\circ u.
\]
Restriction to generators commutes both with precomposition by these
free maps and with postcomposition by any model transformation.
Hence the hom-set bijections are natural in \(X\) and \(M\), proving
\(\mathcal F_{\mathbb T}\dashv U\).

Write \(\mathcal F=\mathcal F_{\mathbb T}\).
For completeness, its counit at \(M\) is the transformation
corresponding to \(1_{M(1)}\); its underlying extension evaluates
each term whose generators are the elements of \(M(1)\).
It satisfies \(U(\varepsilon_M)\circ\eta_{U(M)}=1_{U(M)}\).
The other triangle follows because
\(\varepsilon_{\mathcal F(X)}\circ\mathcal F(\eta_X)\) and the identity
agree on every generator of \(\mathcal F(X)\); freeness makes them equal.
The finite representable theorem and this universal property give
the unique generator-preserving isomorphism
\(\mathcal F_{\mathbb T}(n)\cong R_n\), including \(n=0\).
\end{proof}
\end{annexitem}

\begin{annexitem}{annex:p10.lawvere-free.3}{thm:p10-lawvere-empty-carriers}{The zero-generator and empty-carrier cases}
\begin{proof}
\paragraph{The abstract theory criterion.}
The object \(0\) is terminal in \(\mathbb T\).
Product preservation makes \(c_0^M:M(0)\to M(1)^0\) a bijection,
so \(M(0)\) is a singleton.
If \(t:0\to1\) exists, its image would be a function from this
singleton to \(M(1)\); hence \(M(1)\) cannot be empty.

Conversely suppose \(\mathbb T(0,1)=\varnothing\).
There is no arrow \(0\to n\) for \(n>0\), since composing with
\(p_0^n:n\to1\) would produce one.
Define \(M(0)=\{*\}\) and \(M(n)=\varnothing\) for \(n>0\).
All required functions on arrows now exist uniquely: they are
singleton-to-singleton, empty-to-singleton, or empty-to-empty maps.
Identity and composition laws hold by uniqueness.

The chosen product of finitely many objects \(n_j\) is their sum.
If some \(n_j>0\), its model value and the product of the factor
values are both empty. If all \(n_j=0\), both are singletons.
This includes the empty product. The canonical comparisons are
therefore bijections, proving product preservation.
For any empty-carrier model, its comparisons force singleton value
at \(0\) and empty values at every positive object.
The unique component bijections between any two such models are
natural, because all functions involved are unique.

\paragraph{Ground terms and the initial model.}
In \(T_{\Sigma,E}\), an arrow \(0\to1\) is by definition
a singleton tuple of a ground-term class. A nullary symbol is a ground term.
If no nullary symbol exists, induction on term formation shows that
every term contains a variable; the empty context therefore has no
terms. Passing to a quotient of this empty set cannot add an arrow.
This absence of nullaries also shows that the empty algebra
has all its positive-arity operations and satisfies every equation.

The initial model \(R_0\) has carrier \(\mathbb T(0,1)\), by the
finite free theorem. If this is nonempty, every model has nonempty
carrier by the first paragraph; the two model categories coincide.
If it is empty, \(R_0\) is outside the nonempty full subcategory.
For any nonempty generator set, the term set and its quotient are
nonempty, so its free model does belong to that subcategory.
Its universal property restricts to the nonempty target models.

\paragraph{Nonempty sets and semigroups.}
In the empty signature with no equations, any nonempty set admits
the two distinct constant functions into \(\{0,1\}\).
All functions are homomorphisms, so no such set can be initial.
For semigroups use the two-element algebra \(L=\{0,1\}\) with
\(a b=a\). It is associative and each element is idempotent.
For either \(e\in L\), the constant function from a nonempty
semigroup to \(e\) preserves products:
its value on a product is \(e=e e\).
The two functions are distinct, excluding an initial nonempty semigroup.

\paragraph{A positive-arity operation defining a distinguished value.}
Let \(\Sigma\) contain only a unary \(d\), with equation
\(d(x)=d(y)\) in a two-variable context.
In every nonempty algebra \(A\), the image of \(d\) consists of one
element \(a_*\). Moreover \(d(a_*)=a_*\), since every value of
\(d\) is \(a_*\).
The one-element algebra \(I=\{e\}\), \(d(e)=e\), has a homomorphism
to \(A\) sending \(e\) to \(a_*\).
Any such homomorphism must satisfy
\(h(e)=h(d(e))=d(h(e))=a_*\), so it is unique.
Thus \(I\) is initial in the nonempty subcategory.
There is still no ground term, because the signature has no nullary
symbol. The empty algebra has its unique unary operation and satisfies
the equation vacuously; its unique maps to all algebras preserve \(d\).
It is the all-carrier initial algebra and corresponds to the
empty-carrier initial model.  These are distinct initial-object
statements in two different categories.
\end{proof}
\end{annexitem}
